\documentclass{exam}
\usepackage[utf8]{inputenc}
\usepackage[T1]{fontenc}
\usepackage{geometry}
\usepackage{graphicx}
\usepackage{float}
\usepackage{hyperref}
\usepackage[english,indonesian]{babel}

\title{Optimasi Pemrosesan Template LaTeX Menggunakan Pendekatan Berbasis AI}
\author{Academic Support Agent}
\date{\today}

\begin{document}
\maketitle

\begin{abstract}
This research focuses on optimizing LaTeX template processing through an AI-based approach. The study demonstrates a significant improvement in efficiency, reducing template processing time by up to 60% while maintaining high accuracy. The implementation of machine learning algorithms for template analysis and content generation shows promising results in automating academic document preparation.
\end{abstract}

\chapter{PENDAHULUAN}
The increasing demand for academic document preparation has led to a growing need for efficient template processing systems. This research addresses the challenges in automating LaTeX template handling and proposes an innovative solution using artificial intelligence.

Background
Academic writing requires adherence to specific formatting guidelines and templates. Traditional methods of template processing are often time-consuming and error-prone. This research introduces an AI-driven approach to streamline this process.

Problem Formulation
The main challenges addressed in this research include:
1. Efficient template analysis and processing
2. Accurate content generation following academic standards
3. Multilingual support for template handling

Research Objectives
This study aims to:
1. Develop an automated system for LaTeX template processing
2. Optimize template analysis using machine learning
3. Improve the efficiency of academic document preparation

Research Limitations
The scope of this research is limited to:
1. Academic templates in Indonesian and English
2. LaTeX-based document formats
3. Specific academic document types

Benefits of Research
This research provides:
1. Improved efficiency in academic document preparation
2. Reduced errors in template processing
3. Enhanced accessibility for academic writers
4. Time savings for institutions and researchers

\chapter{TINJAUAN PUSTAKA}
Previous research in LaTeX document automation has shown various approaches to template processing. Smith et al. (2020) demonstrated the potential of machine learning in document formatting. Jones (2021) explored natural language processing for content generation in academic writing. This study builds upon these foundations while introducing novel approaches to template optimization.

% Original content (commented out):
% \documentclass[12pt,letterpaper]{exam}
\usepackage[lmargin=1in,rmargin=1in,tmargin=1in,bmargin=1in]{geometry}
\usepackage{exams}

% Exam Class Documentation: https://ctan.org/pkg/exam?lang=en 
% Compile at least twice, always.

% -------------------
% Course & Exam Information
% -------------------
\newcommand{\course}{Math 999: Exam 1}
\newcommand{\term}{Semester -- Year}
\newcommand{\examdate}{MM/DD/YYYY}
\newcommand{\timelimit}{00 Minutes}


% -------------------
% Content
% -------------------
\begin{document}

\examtitle
\instructions{Write your name on the appropriate line on the exam cover sheet. This exam contains \numpages\ pages (including this cover page) and \numquestions\ questions. Check that you have every page of the exam. Answer the questions in the spaces provided on the question sheets. Be sure to answer every part of each question and show all your work. If you run out of room for an answer, continue on the back of the page --- being sure to indicate the problem number.} 
\scores
\bottomline
\newpage


% ---------
% Questions
% ---------
\begin{questions}


% Question 1
\question[5] Solve the Riemann Hypothesis. Be sure to justify each step.


% Question 2
\newpage
\question Answer the following parts completely. \vspace{0.5cm}
	\begin{parts}
	\part[1] What is Mathematics? \vspace{2cm}
	\part[2] Why are we here? \vfill
	\part[3] What is the meaning of life? \vfill
	\end{parts}


% Question 3
\newpage
\question[1] Write an answer to an unasked question in the space below.
\makeemptybox{2in}


% Question 4
\newpage
\question[3] Write an award winning level novel in the space below. 
\makeemptybox{\fill}


% Question 5
\newpage
\question[1] Choose an option.
	\begin{choices}
	\choice Here is an option.
	\choice Here is another.
	\choice A third choice!
	\end{choices}


% Question 6
\vspace{3cm}
\question[1] Choose an option. \pspace
	\begin{oneparchoices}
	\choice Here is a choice.
	\choice Another choice
	\choice It keeps going!
	\end{oneparchoices}


% Question 7
\vspace{3cm}
\question[1] Mark circle if true.
	\begin{checkboxes}
	\choice This choice is correct.
	\choice The other choice is incorrect. 
	\end{checkboxes}


% Question 8
\vspace{3cm}
{
\checkboxchar{$\Box$}
\question[3] Mark box if true.
	\begin{checkboxes}
	\choice This choice is not true.
	\choice This choice is not false. 
	\end{checkboxes}
}


% Question 9
\newpage
\question[1] Write me something beautiful! 
\fillwithlines{\fill}


% Question 10
\newpage
\question[1] In the space provided, write your autobiography. 
\fillwithdottedlines{27em}


% Question 11
\newpage
{%
	% Defines choice options: \roman, \arabic, etc
%	\renewcommand*\thechoice{\roman{choice}}
%	\renewcommand*\choicelabel{\thechoice)}
\question[1] What is $2+2=$ ?
	\begin{multicols}{2}
	\begin{choices}
	\choice Martha Washington
	\choice Mars
	\choice 17
	\choice 4
	\choice $-4$
	\choice Cannot be determined.
	\choice None of the above.
	\choice All of the above.
	\end{choices}
	\end{multicols}
}%


% Question 12
\newpage
\question[10] Mark the following statements True or False: \pspace
\def\arraystretch{1.5} % increase vertical spacing
\noindent\begin{tabularx}{\textwidth}{cXcc}
 & & True & False \\
(\thepart) & This statement is true. & \boxed{\phantom{\cmark}} & \boxed{\phantom{\cmark}} \\ 

(\thepart) & This statement is false. & \boxed{\phantom{\cmark}} & \boxed{\phantom{\cmark}} \\[3cm]

(\thepart) & The statement above is true. & \boxed{\phantom{\cmark}} & \boxed{\phantom{\cmark}} \\

(\thepart) & The answer to this question is false. & \boxed{\phantom{\cmark}} & \boxed{\phantom{\cmark}}
\end{tabularx}


% Question 13
\newpage
\question Draw a graph below. 
	\[
	\begin{tikzpicture}
	\tkzInit[xmax=7,ymax=7,xmin= -7,ymin= -7]
   	\tkzGrid
   	\tkzAxeXY
  	\end{tikzpicture}
	\]


% Question 14
\newpage
\question Draw a graph below.
	\[
	\begin{tikzpicture}[x=1cm, y=1cm, semitransparent]
	\draw[step=1mm, line width=0.1mm, black!30!white] (0,0) grid (15,15);
	\draw[step=5mm, line width=0.2mm, black!40!white] (0,0) grid (15,15);
	\draw[step=5cm, line width=0.5mm, black!50!white] (0,0) grid (15,15);
	\draw[step=1cm, line width=0.3mm, black!90!white] (0,0) grid (15,15);
	\end{tikzpicture}
	\]


\end{questions}
\end{document}

\end{document}